% handout_template.tex
\documentclass[11pt]{article}
\usepackage{mae2250-handout}

% --- Handout metadata (edit per handout) ---
\renewcommand{\HandoutType}{Assignment} % e.g., Assignment, Open Design Project
\renewcommand{\HandoutNumber}{3} % e.g., 1, 2, ...
\renewcommand{\HandoutTitle}{Parameter-Driven Design}
\renewcommand{\DueDate}{Feb 23 @ 11:59PM}
\renewcommand{\TeamType}{Individual} % e.g., INDIVIDUAL, DESIGN TEAM (4--6 students)
% \renewcommand{\TotalPoints}{15}

\begin{document}
\MakeHandoutHeader

% =========================
% Edit content below
% =========================

\Section{Logistics}

To ensure your deliverables are submitted correctly, and TAs are grading the correct version, you'll need to create a few specific folders in your fusion 2250 project.
First, you'll create a \texttt{<CoffeeCup>} folder in your 2250 project folder.
Then add two sub-folders: one called \texttt{<CoffeeCup\_Working>} and the other labelled \texttt{<CoffeeCup\_Final>}.
As you might expect, you'll create and edit your parts in the working folder.
This is your sandbox!
You can make as many versions, mistakes, or edits as you want here (including accidental destruction and failed repairs).
Your TAs won't grade or look at this folder.
When you are confident that your files are ready to submit, move or copy them to the “Final” folder.
Your TAs will only grade files in this folder – if you do not move them from Working to Final they will not be graded.
This is a good habit to practice not only for this class, but for organizing versions of your files as they go into “Design Freeze” in practice.

\section{Overview}
The goal of this assignment is to learn and become familiar with parametric design.
You client is a new up-and-coming café, and they want to commemorate their anniversary with a line of branded, re-usable cups, aesthetically themed around the classic disposable coffee cup.
You will design a single parametric CAD model of a reusable coffee cup, and use it to then create different cup sizes, driven by a single parameter.
\subsection*{Design Requirements}
Your model must include both the coffee cup and a matching lid.
The café primarily serves commuters and requires that all cups fit securely within a standard car cup holder.
You may assume cup holder diameter for most cars is at least 3 inches.
The café is considering three cup sizes: small, medium, and large, with target volumes of 8 oz, 12 oz, and 16 oz, respectively.
These volumes are provided as reference values only and are not finalized in the eyes of the café team.
Therefore, rather than creating separate CAD files for each size, you should design one model that can be resized to accommodate any volume by changing a single volume parameter.
\subsection*{Other considerations}
Your final submission model must be fully parametric, such that updating the volume parameter automatically adjusts the geometry of the cup and lid as needed.
To achieve this, you are required to use the \texttt{Parameters} feature in Fusion.
\par
\textbf{Plan with a sketch}: before you start CAD, make a sketch of the cup, highlight the relevant dimensions, and use it to plan out which dimensions will be driving parameters, and how they will depend on each other.
Color code them, and annotate your sketch with the equations you'll use.
It should be clear from your sketch what parameter to change to get a desired effect.
Use this process to plan out your CAD timeline.
\par
\textbf{Unit conversion}: in Fusion the parameters feature does not accept conflicting units even if the expression is correct.
To tackle this you can divide by each unit in every expression. For example: $[ 2 * (3 \text{mm}) ]$ will not work, but $[ 2 * (3 \text{mm} / 1 \text{mm}) ]$will.
\par
\textbf{Good Practices}: make sure you are following the \texttt{Good CAD Practices} as posted under assignment 1 in Canvas.
\section{Deliverables}
TAs will inspect your design on Fusion Hub, and will test your parameter-driven design. In addition, submit a PDF with your drawn sketch clearly showing which parameters drive the design, and (typed up) the parameter names.
% \begin{Deliverables}
%   \item \textbf{CAD} of cup and lid in Fusion hub.
%   \item \textbf{PDF} of your drawn sketch, and the equations used.
% %   \item \textbf{Screenshot} of the master sketch, highlighting which parameter to change to drive the overall design.
% \end{Deliverables}

% Example points table rows:
\PointsTableAuto{
  \PointsRow{All CAD sketches, bodies, and functional parameters have a meaningful name}{3}
  \PointsRow{All CAD sketches are fully constrained}{1}
  \PointsRow{All parts are assigned a material, and non-driven dimensions are reasonable for that material}{2}
  \PointsRow{The drawn sketch/equations clearly convey design intent and matches the CAD history}{4}
  \PointsRow{Cup would fit in a standard car cup-holder}{1}
  \PointsRow{Cup capacity can be changed with a single parameter without breaking any part}{4}
}

\end{document}
